\chapter{METHODOLOGY}

\section{Overview}

The analytical framework comprises six modules, applied sequentially after each daily
data collection cycle. Figure~\ref{fig:pipeline} illustrates the processing pipeline.

\begin{figure}[H]
  \centering
  \fbox{\parbox{0.85\textwidth}{\centering
    [Figure: Analytics Pipeline]\\
    Daily Snapshots $\rightarrow$ Entrant/Exit $\rightarrow$ Changes $\rightarrow$
    Price Tier (K-Means) $\rightarrow$ LQS $\rightarrow$ BMS $\rightarrow$ Alerts
  }}
  \caption{Analytics processing pipeline}
  \label{fig:pipeline}
\end{figure}

\section{Entrant and Exit Detection}

New market entrants and products exiting the Top 50 are detected by comparing the
set of ranked ASINs between consecutive daily snapshots:

\begin{align}
  \text{Entrants}_t &= \text{Rankings}_t \setminus \text{Rankings}_{t-1} \\
  \text{Exits}_t    &= \text{Rankings}_{t-1} \setminus \text{Rankings}_t
\end{align}

Events are flagged as \textit{Top 10} entries or exits when the associated rank
falls within the top 10 positions, as these represent commercially significant
movements with greater impact on competing products' visibility.

\section{Listing Quality Score (LQS)}

The Listing Quality Score is a composite metric measuring the completeness and
richness of an Amazon product listing. Inspired by Amazon's own content quality
guidelines \cite{amazon2024listing}, the LQS is computed as a weighted sum of
seven sub-components, each scored on a normalized scale:

\begin{equation}
  \text{LQS} = w_1 S_{\text{title}} + w_2 S_{\text{bullets}} + w_3 S_{\text{image}}
             + w_4 S_{\text{aplus}} + w_5 S_{\text{rating}} + w_6 S_{\text{reviews}}
             + w_7 S_{\text{sentiment}}
\end{equation}

Table~\ref{tab:lqs} defines each sub-component and its weight:

\begin{table}[H]
  \centering
  \begin{tabular}{llcp{5cm}}
    \toprule
    \textbf{Component} & \textbf{Source Field} & \textbf{Max Score} & \textbf{Scoring Rule} \\
    \midrule
    Title completeness  & \texttt{product\_name}   & 10 & 1 pt/word, max 10 \\
    Bullet points       & \texttt{bullet\_count}   & 15 & $\min(n/5 \times 15, 15)$ \\
    Main image          & \texttt{image\_url}      & 15 & 15 if present, else 0 \\
    A+ content          & \texttt{ebc\_html\_hash} & 15 & 15 if present, else 0 \\
    Star rating         & \texttt{stars}           & 20 & $(s/5) \times 20$ \\
    Review volume       & \texttt{reviews\_count}  & 15 & $\min(n/1000 \times 15, 15)$ \\
    Review sentiment    & \texttt{avg\_sentiment}  & 10 & $(\hat{s}+1)/2 \times 10$ \\
    \midrule
    \textbf{Total}      &                          & \textbf{100} & \\
    \bottomrule
  \end{tabular}
  \caption{LQS component definitions and weights}
  \label{tab:lqs}
\end{table}

\section{Brand Momentum Score (BMS)}

The Brand Momentum Score captures a product's competitive trajectory by combining
three velocity signals:

\begin{equation}
  \text{BMS}_t = 0.5 \cdot \hat{v}_{\text{BSR}} + 0.3 \cdot \hat{v}_{\text{reviews}}
               + 0.2 \cdot \hat{s}_{\text{sentiment}}
\end{equation}

where each component is normalized to $[0, 1]$:

\begin{align}
  v_{\text{BSR},t}     &= \frac{\text{BSR}_{t-1} - \text{BSR}_t}{\text{BSR}_{t-1}}
  \quad \text{(positive = improving rank)} \\
  v_{\text{reviews},t} &= \text{ReviewCount}_t - \text{ReviewCount}_{t-1} \\
  s_{\text{sentiment}} &= \text{avg\_sentiment\_score} \in [-1, 1]
\end{align}

BSR velocity receives the highest weight (0.5) as it most directly reflects current
sales momentum. Review velocity (0.3) captures medium-term organic growth. Sentiment
(0.2) provides a leading indicator of future rating changes.

\section{K-Means Price Tiering}

Products within each category are segmented into three price tiers — \textit{entry},
\textit{mid-range}, and \textit{premium} — using K-Means clustering applied
independently per category per day.

Given a set of product prices $\{p_1, p_2, \ldots, p_n\}$ within a category,
K-Means minimizes the within-cluster sum of squares:

\begin{equation}
  \underset{C}{\arg\min} \sum_{k=1}^{K} \sum_{p_i \in C_k} (p_i - \mu_k)^2
\end{equation}

Prices are standardized using $z$-score normalization before clustering to ensure
numerical stability. Cluster labels are assigned post-hoc by ranking cluster centroids:
the lowest-centroid cluster is labeled \textit{entry}, the middle \textit{mid-range},
and the highest \textit{premium}.

The algorithm is initialized with $k=3$ and $n\_init=10$ random seeds to ensure
convergence to a stable solution. When fewer than three distinct price points exist in
a category (rare), $k$ is reduced accordingly.

\section{Sentiment Analysis with RoBERTa}

\subsection{Model Selection}

We employ the \texttt{cardiffnlp/twitter-roberta-base-sentiment-latest} model, a
RoBERTa-base architecture fine-tuned on approximately 124 million tweets and subsequently
fine-tuned on a sentiment classification task \cite{loureiro2022timelms}. Despite being
trained on social media text, this model has demonstrated strong transfer to product
review domains due to the informal, opinionated nature of both text types.

\subsection{Inference Pipeline}

For each review, the \texttt{reviewText} field is truncated to 512 tokens (the model's
maximum context length) and passed to the classification head, producing a probability
distribution over three classes: \{positive, neutral, negative\}.

A scalar sentiment score $s \in [-1, 1]$ is derived as:
\begin{equation}
  s = \text{confidence} \times
  \begin{cases}
    +1.0 & \text{if label = positive} \\
     0.0 & \text{if label = neutral} \\
    -1.0 & \text{if label = negative}
  \end{cases}
\end{equation}

\subsection{Aggregation}

Per-review scores are aggregated to the ASIN level:

\begin{align}
  \bar{s}_{\text{ASIN}} &= \frac{1}{N} \sum_{i=1}^{N} s_i \\
  r_{\text{pos}}        &= \frac{|\{i : \text{label}_i = \text{positive}\}|}{N} \\
  r_{\text{neg}}        &= \frac{|\{i : \text{label}_i = \text{negative}\}|}{N}
\end{align}

\subsection{Aspect-Level Analysis}

The \texttt{web\_wanderer} actor additionally provides pre-computed aspect-level
sentiment annotations, identifying sentiment polarity for specific product attributes
(e.g., ``Quality'', ``Battery Life'', ``Sound'', ``Value''). These aspects serve as a
complementary signal and validation benchmark for the RoBERTa-based classification.

\section{Price Trend Forecasting with Prophet}

\subsection{Model Specification}

Facebook Prophet \cite{taylor2018forecasting} decomposes price time series as:

\begin{equation}
  y(t) = g(t) + s(t) + h(t) + \epsilon_t
\end{equation}

where $g(t)$ is the trend component, $s(t)$ captures seasonality (disabled given the
short observation window), $h(t)$ represents holiday effects (not applicable), and
$\epsilon_t$ is the error term.

Given the 15-day collection window, we fit a linear trend model ($g(t) = k + at$)
with no seasonality components. The model produces a 7-day ahead forecast with 80\%
uncertainty intervals.

\subsection{Limitations Acknowledgement}

It is important to note that price forecasting with $n \approx 15$ observations is
inherently limited in predictive accuracy. The primary value of the forecasting module
in this study is \textbf{methodological demonstration} — establishing the pipeline
for future deployment with longer data histories — rather than producing operationally
reliable price predictions.

\section{Alert Engine}

A rule-based alert engine monitors daily snapshots and derived metrics to generate
actionable notifications. Table~\ref{tab:alerts} defines the alert rules.

\begin{table}[H]
  \centering
  \begin{tabular}{llll}
    \toprule
    \textbf{Alert Type} & \textbf{Trigger Condition} & \textbf{Severity} & \textbf{Source} \\
    \midrule
    \texttt{price\_drop}   & Price decrease $\geq 15\%$ day-over-day  & Medium/High & daily\_snapshots \\
    \texttt{stockout}      & \texttt{in\_stock} flips False              & High        & daily\_snapshots \\
    \texttt{bsr\_drop}     & BSR worsens by $\geq 500$ positions        & Medium      & daily\_snapshots \\
    \texttt{image\_change} & pHash Hamming distance $> 10$              & Low         & image\_change\_events \\
    \texttt{new\_entrant}  & ASIN appears in rankings for first time    & Low/High    & entrant\_exit\_events \\
    \texttt{exit}          & ASIN disappears from Top 50                & Medium      & entrant\_exit\_events \\
    \bottomrule
  \end{tabular}
  \caption{Alert engine rule definitions}
  \label{tab:alerts}
\end{table}

\section{Multi-Agent Delivery Layer}

The analytical outputs above are stored in Supabase but only become useful when
they reach the decision-maker. Rather than relying solely on the Streamlit
dashboard, this thesis adopts a \textbf{multi-agent delivery architecture}
built on top of \textbf{OpenClaw}, a self-hosted agent runtime that connects
large language models (LLMs) to user-defined tools (``skills'') over messaging
channels.

\subsection{Architecture}

Three specialist agents are deployed, each bound 1-to-1 to a dedicated
Telegram bot account. The binding is performed by the OpenClaw Gateway, which
multiplexes messages between channels, agents, and skills. Figure~\ref{fig:openclaw}
summarises the topology.

\begin{figure}[H]
  \centering
  \fbox{\parbox{0.92\textwidth}{\centering
    [Figure: OpenClaw multi-agent delivery]\\[2pt]
    \texttt{@babyspyyy\_bot} $\to$ Competitor Spy (6 skills)\\
    \texttt{@babydetective\_bot} $\to$ Sentiment Detective (4 skills)\\
    \texttt{@babystrategist\_bot} $\to$ Momentum Strategist (5 skills)\\[2pt]
    Gateway $\to$ Agents $\to$ Python skill CLI $\to$ Supabase
  }}
  \caption{OpenClaw multi-agent delivery architecture}
  \label{fig:openclaw}
\end{figure}

The split mirrors the analytical structure of the pipeline:

\begin{itemize}
  \item \textbf{Competitor Spy} surfaces competitor-side signals: BMS,
  category rankings, sponsored slot share, entrant/exit events, price-tier
  shifts and image redesigns.
  \item \textbf{Sentiment Detective} surfaces customer-side signals: per-ASIN
  RoBERTa sentiment trends, raw quote samples, and aspect-level breakdowns.
  \item \textbf{Momentum Strategist} synthesises the two views into an
  actionable daily brief covering alerts, BMS top movers, LQS, and the 7-day
  price forecast.
\end{itemize}

\subsection{Skill Contract}

Every analytical capability exposed to an agent is implemented as a Python
\textbf{skill}: a command-line program that reads JSON arguments from
\texttt{stdin} and returns JSON-serialisable results on \texttt{stdout}. A
machine-readable schema (\texttt{python skill.py --schema}) declares the
skill's name, parameters, and description. Thirteen such skills (Table~\ref{tab:skills})
are wired into OpenClaw via per-skill YAML manifests under
\texttt{openclaw/wrappers/}. All skills are \textbf{read-only}: they query
Supabase but never write back. Mutations remain confined to the offline
ingest and analytics scripts described in earlier sections.

\begin{table}[H]
  \centering
  \begin{tabular}{lll}
    \toprule
    \textbf{Group} & \textbf{Skill} & \textbf{Purpose} \\
    \midrule
    sentiment & \texttt{query\_sentiment}      & Per-ASIN sentiment time series \\
    sentiment & \texttt{query\_reviews}        & Recent raw review quotes \\
    sentiment & \texttt{query\_aspects}        & Aspect-level sentiment summary \\
    market    & \texttt{query\_bms}            & Top BMS movers per category \\
    market    & \texttt{query\_rankings}       & Top-50 category snapshot \\
    market    & \texttt{query\_entrant\_exits} & Entrant / exit events \\
    market    & \texttt{query\_sponsored\_share} & Share-of-voice in sponsored slots \\
    market    & \texttt{query\_price\_tiers}   & K-Means price tier membership \\
    market    & \texttt{query\_price\_forecast} & 7-day Prophet forecast \\
    listing   & \texttt{query\_snapshots}      & Per-ASIN daily snapshot detail \\
    listing   & \texttt{query\_image\_changes} & pHash change events \\
    listing   & \texttt{query\_lqs}            & LQS score and sub-components \\
    alerts    & \texttt{query\_alerts}         & Active alerts by severity \\
    \bottomrule
  \end{tabular}
  \caption{The 13 OpenClaw skills exposed to the agent layer}
  \label{tab:skills}
\end{table}

\subsection{Agent Voice and Output Contract}

Each agent has its own workspace under \texttt{openclaw/agents/<name>/}
containing four markdown files: \texttt{SOUL.md} (voice and rules),
\texttt{IDENTITY.md} (display name, emoji, vibe), \texttt{AGENTS.md} (skill
list with JSON examples and trigger keywords in English and Vietnamese), and
\texttt{USER.md} (user context). The \texttt{SOUL.md} files specify a
production-hardened output contract: hard line caps per message type,
no markdown tables (which break on Telegram mobile), backtick-wrapped
numbers and ASINs, and a banned set of conversational closers
(``Let me know if\ldots'', ``Hope this helps'') so that every reply ends in a
concrete next-step or a low-confidence caveat.

\subsection{Scheduled Briefs}

In addition to interactive queries, three cron jobs (managed by the OpenClaw
Gateway, time zone \texttt{Asia/Ho\_Chi\_Minh}) push proactive briefs:

\begin{itemize}
  \item \texttt{daily-brief-strategist} --- 08:00 daily; the Strategist
  produces a three-list summary (alerts, BMS top movers, LQS bottom).
  \item \texttt{weekly-competitor-spy} --- 09:00 Monday; the Spy summarises
  the previous week's competitor moves.
  \item \texttt{weekly-sentiment-detective} --- 09:00 Wednesday; the
  Detective digests the previous week's review sentiment shifts.
\end{itemize}

This separation between (a) interactive question-answering and (b) scheduled
push briefs mirrors the dual mode in which decision-makers consume
intelligence in practice.

\section{Evaluation Methodology}

Two of the analytical modules in this thesis --- the RoBERTa sentiment
classifier and the Prophet price forecaster --- produce predictions whose
quality must be quantified before the downstream BMS, LQS, and forecast
skills can be trusted. This section formalises the two evaluation protocols.

\subsection{Sentiment Classification --- Distant Supervision against Stars}

A gold-labelled in-domain sentiment dataset is not available within the scope
of this thesis. We therefore adopt a \textbf{distant-supervision} protocol
\cite{go2009twitter}: the reviewer's numeric star rating is treated as a
noisy ground-truth signal for the underlying sentiment polarity. The mapping
follows the 3-class convention used in the literature for Amazon and Yelp
reviews \cite{mcauley2013amateurs}:

\begin{equation}
  y_{\text{star}} = \begin{cases}
    \text{negative} & \text{if rating} \in \{1, 2\} \\
    \text{neutral}  & \text{if rating} = 3 \\
    \text{positive} & \text{if rating} \in \{4, 5\}
  \end{cases}
\end{equation}

For each labelled review, the RoBERTa pipeline produces a predicted class
$\hat{y}_{\text{RoBERTa}}$. We report \textbf{accuracy}, \textbf{macro-averaged
F\textsubscript{1}}, per-class precision/recall/F\textsubscript{1}, and the
full $3\times3$ confusion matrix. The macro-F\textsubscript{1} is computed as

\begin{equation}
  \text{macro-F}_1 = \frac{1}{3} \sum_{c \in \{\text{neg, neu, pos}\}}
    \frac{2 \cdot P_c \cdot R_c}{P_c + R_c}
\end{equation}

so that the heavy class-imbalance towards positive reviews does not mask
poor performance on the smaller negative and neutral classes.

\subsubsection{Limitations of Distant Supervision}

The star-rating mapping is a noisy proxy for two well-documented reasons:
(a) the relationship between numeric rating and free-text sentiment is not
deterministic (a 4-star review can contain substantial criticism, and a
1-star review can be neutral in tone if the complaint concerns shipping
rather than the product); and (b) the 3 / neutral bucket is particularly
sparse and ambiguous, since reviewers tend to gravitate towards the
extremes \cite{mcauley2013amateurs}. The reported accuracy and F\textsubscript{1}
should therefore be read as lower bounds on the underlying classifier
quality. Despite this, distant supervision provides a defensible benchmark
and is consistent with practice in the e-commerce sentiment literature where
gold-labelled review data is unavailable.

\subsection{Price Forecasting --- Walk-Forward Backtest}

For each watchlist ASIN with at least \textbf{seven} days of price history,
we evaluate the Prophet forecaster using a \textbf{walk-forward expanding
window} backtest:

\begin{enumerate}
  \item Reserve the most recent $H = 3$ days as a hold-out test set.
  \item Train Prophet on the remaining (oldest) $n - H$ days.
  \item Produce a 3-day-ahead forecast with 80\% prediction interval (PI).
  \item Compute prediction error against the actual prices in the hold-out.
\end{enumerate}

This protocol respects temporal ordering --- the forecaster never sees
future prices --- and matches the production deployment scenario in which
Prophet is retrained daily on all historical data and used to predict the
next 7 days.

\subsubsection{Metrics}

For each ASIN $i$ with hold-out actuals $\{y_{i,t}\}$ and forecast means
$\{\hat{y}_{i,t}\}$, we compute:

\begin{align}
  \text{MAE}_i  &= \tfrac{1}{H}\sum_{t} |y_{i,t} - \hat{y}_{i,t}| \\
  \text{MAPE}_i &= \tfrac{1}{H}\sum_{t} \frac{|y_{i,t} - \hat{y}_{i,t}|}{y_{i,t}} \times 100\%
\end{align}

The pooled MAE/MAPE across all evaluated ASINs are reported in
Chapter~\ref{chap:results}. To verify that Prophet adds value over the
simplest possible baseline, we additionally compute a \textbf{naive last-value
baseline} ($\hat{y}_{i,t} = y_{i, T-H}$ for all $t > T-H$, i.e.\ assume the
price stays flat at its most recent training value) and report its MAE
alongside the Prophet MAE. Finally, we compute \textbf{prediction interval
coverage} --- the share of hold-out points falling inside Prophet's 80\% PI
--- as a diagnostic of probabilistic calibration.
